\documentclass{article}

\usepackage[english]{babel}
\usepackage[letterpaper,top=2cm,bottom=2cm,left=3cm,right=3cm]{geometry}
\usepackage{amsmath,amssymb}
\setlength{\parindent}{0pt}

\title{DSC 190/291: Learning Theory through Formal Proof and Proof Presentation\\Assignment 3}
\date{UCSD, Spring 2026}

\begin{document}
\maketitle

\section{Overview}

This assignment has one large technical mini-course, one theorem/comparison part, and an AI usage report:
\begin{itemize}
\item Part A: a mini-course on concentration inequalities.
\item Part B: the No-Free-Lunch theorem and the Fundamental Theorem of PAC learning.
\item Part C: AI usage report.
\end{itemize}

\section{Part A: A Mini-Course on Concentration Inequalities (65 points)}

In the Week 3 lecture we proved the following i.i.d. growth-function ERM guarantee. Let
$$
\widehat{h} \in \arg\min_{h \in \mathcal{H}} L_n(h).
$$
Then with high probability over $S \sim \mathcal{D}^n$,
$$
L_{\mathcal{D}}(\widehat{h})
\le
\inf_{h \in \mathcal{H}} L_{\mathcal{D}}(h)
+ O\!\left(
\sqrt{\frac{\log \Gamma_{\mathcal{H}}(2n) + \log(1/\delta)}{n}}
\right).
$$

The learning-theoretic structure of the proof was covered in lecture, but several concentration inequalities were used without proof. The inequalities to learn in this mini-course are:
\begin{itemize}
\item Hoeffding's inequality,
\item Hoeffding's inequality for sampling without replacement,
\item McDiarmid's inequality,
\item Bernstein's inequality.
\end{itemize}

Create a short mini-course that teaches these four concentration inequalities well enough to complete the Week 3 ERM proof end-to-end. The mini-course should be one coherent document written so that someone unfamiliar with the material could learn from it. Organize it in roughly the following order:
\begin{enumerate}
\item The concept of concentration.
\item The common technique.
\item Statements, assumptions, and consequences.
\item Proofs.
\item Connection to the Week 3 ERM proof.
\end{enumerate}

Provide detailed references for all results you used, and make sure everything can be checked against those references.

\section{Part B: The No-Free-Lunch Theorem and the Fundamental Theorem (20 points)}

In this part you will prove the i.i.d. No-Free-Lunch theorem, use it to prove the lower-bound direction of the Fundamental Theorem of PAC learning, construct an explicit bad example for a learning rule of your choice, and compare the Week 1 and Week 3 NFL theorems.

\section{Part C: AI Usage Report (15 points)}

Write a short report describing how you used AI in this assignment, including concrete accepted suggestions, rejected suggestions, verification steps, and the five most recent changes you made to your AI workflow.

\end{document}
